\clearpage
\hypertarget{native-482-june4}{}
\begin{center}
\rule{0.55\linewidth}{0.4pt}\\[2pt]
Astronomische Nachrichten. Band XXVII. Nr. 641. S. 265. 266. 29. Juni 1848.\\[2pt]
\rule{0.55\linewidth}{0.4pt}
\end{center}
\begin{flushright}Göttingen, 4. Juni 1848.\end{flushright}

Zu den beiden Beobachtungen von \textsc{Graham}s Planeten, die ich Ihnen am 22. Mai geschickt habe, kann ich nur noch eine weitere hinzufügen. Auch sie wurde am Meridiankreis gemacht:

\begin{center}
1848 Mai 24.\quad $10^{\rm u}\ 18^{\rm m}\ 39^{\rm s}0$ M. Z.\\
Rektasz. $217^\circ\ 15'\ 58''8$\qquad Dekl. $11^\circ\ 29'\ 29''65$ südl.
\end{center}

Der Planet war fast noch schwerer zu beobachten als am 11. Da sich der Zeitraum nun aber fast verdoppelt hatte, veranlasste ich Professor \textsc{Goldschmidt} gleich danach, die Bahn neu zu berechnen. Dabei ergaben sich die folgenden Elemente:

\begin{center}
\textsc{Graham}s Planet.\\[3pt]
\small
\setlength{\tabcolsep}{3pt}
\begin{tabular}{lr}
Epoche, mittl. Länge 1848 Mai 1. $0^{\rm u}$ M. Z. Berl. \mbox{.\ .} & $215^\circ\ 23'\ 57''6$ \\
Perihel \mbox{.\ .\ .\ .\ .\ .\ .\ .\ .\ .\ .\ .\ .\ .\ .\ .\ .} & $71\ 41\ 23.9$ \\
Aufsteigender Knoten \mbox{.\ .\ .\ .\ .\ .\ .\ .\ .\ .\ .\ .} & $68\ 29\ 31.5$ \\
\multicolumn{2}{l}{\quad Alle drei vom mittleren Äquinoktium des 1. Mai 1848 aus gerechnet.} \\
Neigung der Bahn \mbox{.\ .\ .\ .\ .\ .\ .\ .\ .\ .\ .\ .\ .} & $5\ 35\ 23.2$ \\
Exzentrizitätswinkel \mbox{.\ .\ .\ .\ .\ .\ .\ .\ .\ .\ .\ .\ .} & $7\ 6\ 17.7$ \\
Logarithmus der großen Halbachse \mbox{.\ .\ .\ .\ .\ .\ .} & $0.3776396$ \\
Mittlere tägliche siderische Bewegung \mbox{.\ .\ .\ .\ .\ .} & $962''8247$ \\
\end{tabular}
\end{center}

Diese Elemente beruhten auf der Beobachtung von \textsc{Graham} vom 26. April, der Berliner Beobachtung vom 11. Mai und meiner eigenen vom 24. Mai. Professor \textsc{Goldschmidt} hat diese Elemente mit allen Meridianbeobachtungen verglichen, die mir bisher bekannt geworden sind. Diesen Vergleich lege ich auf einem gesonderten Blatt bei. Meine Beobachtungen am Meridiankreis sind unterbrochen, weil das Instrument für eine Änderung jetzt abgenommen ist.
